% =====================================================================
% LMI001 -- theory/invariances.tex  (3 October 2026)
% Blocks ready to paste into manuscript/main.tex (v0.3).  Same macros:
% \R, \one, \tr, \status{...}; theorem environments of main.tex.
% Nothing here has been pasted; main.tex is unchanged.  Each block is
% delimited by "%%%% BLOCK X BEGIN/END" (the compile test splits on them).
%
% WHERE EACH BLOCK GOES, AND WHICH SENTENCES OF v0.3 CHANGE
%  BLOCK A  -> Section 3.1, right after the paragraph that follows the
%              proof of Lemma 3.1 ("In particular the minimisers ... below.").
%              Proposition "Lemma 3.1 is complete" + Corollary on
%              representatives.
%  BLOCK B  -> right after the proof of Proposition 3.3 (before Remark 3.4).
%              Corollary on the rest-length spring + Remark on the weaker
%              (same-minimiser) notion.
%              In the "Rest length" paragraph of the proof of Prop. 3.3,
%              REPLACE the sentence
%                "Whether Lemma~\ref{lem:inv} lists all the transformations
%                 that leave $J$ unchanged up to a constant on every dataset,
%                 and hence whether some other dataset-independent PD
%                 representative exists, we leave open."
%              BY
%                "By Proposition~\ref{prop:invcomplete} this class contains
%                 every dataset-independent representative, so none is PD
%                 (Corollary~\ref{cor:rest})."
%  BLOCK C  -> Section 4, after the paragraph "Values, not minimisers." and
%              before "The dynamic route."  Definition of weighted kernel
%              k-means + Propositions on size weights (Next step (c)).
%  BLOCK D  -> replaces Open problem 4.2 (the statement as written is
%              satisfied by a = 0; see Remark) and adds the Remark after it.
%  OTHER SENTENCES TO CHANGE (text suggestions, not a block):
%   * Prop. 3.5, first line: after "(for instance $K+\sigma I$ or, for $\phi$
%     of negative type, $\widetilde K$)" add "; by
%     Corollary~\ref{cor:reps} these are exactly the PSD matrices
%     $K+\sigma I+g\one^{\!\top}+\one g^{\!\top}$".
%   * Table 2 (claims): replace the row
%       "$(d-r)^2$ is not of negative type; no PD kernel ... & proved here
%        (elementary); other representatives open"
%     by
%       "$(d-r)^2$ is not of negative type; it has no dataset-independent
%        PD representative of any kind (Cor.~\ref{cor:rest}) & proved here
%        (elementary)"
%     and add the rows
%       "Lemma 3.1 lists all value invariances of $J$ for $2\le k\le n-1$
%        (Prop.~\ref{prop:invcomplete}) & proved here (elementary); converse
%        of \citealp{roth2003}; not found in the literature we know"
%       "Fixed point weights: $\sum_cw(n_c)\sum A_{ij}$ is weighted kernel
%        $k$-means for all $A$ iff $w\propto1/m$ ($3\le k\le n-2$)
%        (Prop.~\ref{prop:weights}) & proved here"
%       "Total energy with $A\ge0$ is not weighted kernel $k$-means for any
%        positive point weights (Prop.~\ref{prop:weightsdd}) & proved here"
%   * Next steps: delete (c) and (d); keep (a), (b), (e); add
%     "(c') the case $k=2$ of Proposition~\ref{prop:weights} when
%      $h(a)+h(n-a)$ is constant, and data-dependent weights for decreasing
%      $w$ other than $1/m$; (d') whether the rest-length spring has a
%      dataset-independent PD kernel with the same minimisers
%      (Remark~\ref{rem:argmin})."
%   * Contribution paragraph / abstract (optional): "...; a proof that
%     Lemma 3.1 lists all invariances of the trace form, so that the
%     rest-length spring has no dataset-independent PD representative; and
%     a characterisation of the size weights that are weighted kernel
%     $k$-means objectives (only $1/n_c$)".
%   * Limitations: "The normalisation family is $\{1,1/n_c\}$" can add
%     "(Proposition~\ref{prop:weights} shows that $1/n_c$ is the only size
%      weight inside weighted kernel $k$-means with fixed point weights)".
%   * README / CONTINUIDAD: update the corresponding items (not done here).
% Verification: theory/check_invariances.py (exact ranks for n<=7, all k),
% output in theory/check_invariances_output.txt.
% =====================================================================

%%%% BLOCK A BEGIN
\begin{proposition}[Lemma~\ref{lem:inv} is complete]\label{prop:invcomplete}
Let $n\ge3$, $2\le k\le n-1$ and $\Delta\in\R^{n\times n}$ symmetric. Then $J_\Delta(C)$ takes the same value on all partitions of $[n]$ into $k$ non-empty clusters if and only if
\[
\Delta=\sigma I+g\one^{\!\top}+\one g^{\!\top}\qquad\text{for some }\sigma\in\R,\ g\in\R^n,
\]
and then $J_\Delta\equiv\sigma(n-k)$. The pair $(\sigma,g)$ is unique, so this space $V_{n,k}$ has dimension $n+1$ and does not depend on $k$; it is also the space of $\Delta$ with $J_\Delta(C)=\alpha+\beta(n-k)$ for all partitions with any number of clusters (and then $\alpha=0$, $\beta=\sigma$). For $k\in\{1,n\}$ there is a single partition and every $\Delta$ qualifies. Equivalently, for $2\le k\le n-1$: $J_K-J_{K'}$ is constant on $k$-partitions iff the induced dissimilarities $D(K)_{ij}=K_{ii}+K_{jj}-2K_{ij}$ and $D(K')$ differ by a constant off the diagonal, iff $P(K-K')P\in\R P$ with $P=I-\one\one^{\!\top}/n$. \status{proved here (elementary)}
\end{proposition}
\begin{proof}
\emph{If.} Lemma~\ref{lem:inv}, with value $\sigma(n-k)$ (the matrix $a\one\one^{\!\top}$ is the case $g=\tfrac a2\one$).
\emph{Uniqueness.} If $\sigma I+g\one^{\!\top}+\one g^{\!\top}=0$, the off-diagonal entries give $g_i+g_j=0$ for $i\ne j$; with a third index $l$, $g_i=\tfrac12[(g_i+g_j)+(g_i+g_l)-(g_j+g_l)]=0$, and then $\sigma=0$.
\emph{Only if.} Put $f=\tfrac12(\Delta_{11},\dots,\Delta_{nn})$ and $e=\Delta-f\one^{\!\top}-\one f^{\!\top}$, which has zero diagonal and, by Lemma~\ref{lem:inv}, $J_e=J_\Delta$. For zero diagonal, $J_e(C)=-2\,\mathcal E(C)$ with $\mathcal E(C)=\sum_c\frac1{n_c}\sum_{i<j\in C_c}e_{ij}$, so $\mathcal E$ is constant on $k$-partitions. We show $e_{pi}=e_{pj}$ for all distinct $p,i,j$. If $k\ge3$, the set $R=[n]\setminus\{p,i,j\}$ has $n-3\ge k-2\ge1$ elements; fix a partition of $R$ into $k-2$ non-empty blocks and compare the partitions $\{\{p,i\},\{j\},R\text{-blocks}\}$ and $\{\{p,j\},\{i\},R\text{-blocks}\}$: their energies differ by $\tfrac12(e_{pi}-e_{pj})$, which must vanish. If $k=2$, let $U=[n]\setminus\{i,j\}\neq\emptyset$ and $\gamma_x=e_{ix}-e_{jx}$ for $x\in U$. For $P\subseteq U$ and $Q=U\setminus P$, the partitions $\{P\cup\{i\},Q\cup\{j\}\}$ and $\{P\cup\{j\},Q\cup\{i\}\}$ have energies differing by $\frac1{|P|+1}\sum_{x\in P}\gamma_x-\frac1{|Q|+1}\sum_{x\in Q}\gamma_x$. With $P=\emptyset$ this gives $\sum_{x\in U}\gamma_x=0$; with $P=\{p\}$ it gives $\gamma_p\bigl(\tfrac12+\tfrac1{n-2}\bigr)=0$, so $e_{ip}=e_{jp}$. In both cases, any two pairs are linked by pairs sharing an index ($e_{ab}=e_{ad}=e_{cd}$), so $e_{ij}=-\sigma$ for all $i\ne j$ and some $\sigma$, i.e.\ $e=\sigma I-\sigma\one\one^{\!\top}$ and $\Delta=\sigma I+g\one^{\!\top}+\one g^{\!\top}$ with $g=f-\tfrac\sigma2\one$.
\emph{All $k$.} The condition for all partitions implies the one for $k=2$, hence the form; conversely that form gives $\sigma(n-k)$ for every $k$.
\emph{Equivalent forms.} $D$ is linear with kernel $\{g\one^{\!\top}+\one g^{\!\top}\}$ (if $D(K)=0$ then $K_{ij}=\tfrac12(K_{ii}+K_{jj})$), and $D(\sigma I)=2\sigma(\one\one^{\!\top}-I)$; likewise $M\mapsto PMP$ has the same kernel on symmetric matrices and $P(\sigma I)P=\sigma P$.
\end{proof}

\noindent Since $J_K(C)=\sum_c\frac1{n_c}\sum_{i<j\in C_c}D(K)_{ij}$, the proposition says that the pairwise-clustering cost of a dissimilarity matrix~\citep{hofmann1997} determines the off-diagonal dissimilarities up to one additive constant: it is the converse of the shift invariance of \citet{roth2003}. We did not find the converse stated in the literature we know, but have not searched exhaustively. The exact ranks are confirmed for $n\le7$ and all $k$ by \path{theory/check_invariances.py}.

\begin{corollary}[All representatives, and the minimal shift]\label{cor:reps}
Let $K$ be symmetric, $n\ge3$, $2\le k\le n-1$. The symmetric $K'$ with $J_{K'}-J_K$ constant on $k$-partitions are exactly $K'=K+\sigma I+g\one^{\!\top}+\one g^{\!\top}$. Such a $K'$ can be chosen PSD iff $\sigma\ge\sigma_*(K):=\max\{-c^{\!\top}Kc:\ c\perp\one,\ \lVert c\rVert=1\}$; for $\sigma\ge\sigma_*(K)$, $K'=P(K+\sigma I)P$ is one. \status{proved here; in dissimilarity language the threshold is the minimal constant shift of \citet{roth2003}}
\end{corollary}
\begin{proof}
The first sentence is Proposition~\ref{prop:invcomplete} applied to $K'-K$. If $K'\succeq0$ and $c\perp\one$, then $0\le c^{\!\top}K'c=c^{\!\top}Kc+\sigma\lVert c\rVert^2$, so $\sigma\ge\sigma_*(K)$. Conversely $P(K+\sigma I)P=K+\sigma I+g\one^{\!\top}+\one g^{\!\top}$ for some $g$ (expand $P=I-\one\one^{\!\top}/n$ and absorb $\one\one^{\!\top}$ into $g$), and it is PSD iff $c^{\!\top}(K+\sigma I)c\ge0$ for all $c\perp\one$.
\end{proof}
%%%% BLOCK A END

%%%% BLOCK B BEGIN
\begin{corollary}[The rest-length spring has no dataset-independent PD representative]\label{cor:rest}
Let $\phi\colon[0,\infty)\to\R$ and $A[X]_{ij}=\phi(\lVert x_i-x_j\rVert)$. Call a symmetric $\kappa\colon\R^D\times\R^D\to\R$ a \emph{dataset-independent representative} of $\phi$ if for every set $X\subset\R^D$ of $n\ge3$ distinct points and some $k\in\{2,\dots,n-1\}$ (equivalently, by Proposition~\ref{prop:invcomplete}, every such $k$), $J_{\kappa[X]}-J_{-A[X]}$ is constant on the partitions of $X$ into $k$ clusters.
\begin{enumerate}[label=(\alph*),leftmargin=2em]
\item $\kappa$ is a dataset-independent representative iff there are $g\colon\R^D\to\R$ and $\sigma\in\R$ with $\kappa(x,y)=-\phi(\lVert x-y\rVert)+g(x)+g(y)$ for $x\ne y$ and $\kappa(x,x)=-\phi(0)+\sigma+2g(x)$. This is the class $\widetilde K+f(x)+f(y)+a+\sigma[x=y]$ of Proposition~\ref{prop:bernstein}.
\item For $\phi(d)=(d-r)^2$, $r>0$, no dataset-independent representative is PD. Quantitatively, on $X_s=\{0,su,2su,3su\}$ ($\lVert u\rVert=1$) the minimal shift of Corollary~\ref{cor:reps} satisfies $\sigma_*(-A[X_s])\ge2rs$, which is unbounded in $s$.
\item (b) still holds if $J_{\kappa[X]}$ is only required to equal $\lambda_XJ_{-A[X]}$ plus a constant, with a factor $\lambda_X>0$ that may depend on $X$.
\end{enumerate}
\status{proved here (elementary)}
\end{corollary}
\begin{proof}
(a) \emph{If}: Lemma~\ref{lem:inv}. \emph{Only if}: by Proposition~\ref{prop:invcomplete}, for each $X$ there is a unique $(\sigma_X,g_X)$ with $\kappa[X]+A[X]=\sigma_XI+g_X\one^{\!\top}+\one g_X^{\!\top}$. If $X\subseteq Y$, the principal submatrix of $Y$'s decomposition is a decomposition for $X$, so by uniqueness $\sigma_X=\sigma_Y$ and $g_X=g_Y|_X$. Comparing any $X,X'$ through $X\cup X'$ shows that $\sigma_X=\sigma$ is one constant and that $g(x):=g_X(x)$ does not depend on the set $X\ni x$. (b) If $\kappa$ were PD, then with $c=(1,-1,-1,1)$ on $X_s$, the terms in $g$ vanish against $\sum_ic_i=0$, and $\sum_{i,j}c_ic_j\kappa(x_i,x_j)=-\sum_{i,j}c_ic_j\phi(d_{ij})+4\sigma=-8rs+4\sigma$ (the computation of Proposition~\ref{prop:bernstein}), which is negative for $s>\sigma/2r$. Also $c/2$ is a unit vector orthogonal to $\one$ with $(c/2)^{\!\top}A[X_s](c/2)=2rs$, whence $\sigma_*\ge2rs$. (c) Now $\kappa[X]=\lambda_X(-A[X])+\sigma_XI+g_X\one^{\!\top}+\one g_X^{\!\top}$. For $X\subseteq Y$, restriction gives $(\lambda_X-\lambda_Y)A[X]\in V_{|X|,2}$. The matrix $A[X_s]$ is not in $V_{4,2}$, because $D(A[X_s])_{ij}=2\phi(0)-2\phi(d_{ij})$ is not constant off the diagonal ($(s-r)^2,(2s-r)^2,(3s-r)^2$ are never all equal for $s>0$). Hence $\lambda_{X_s}=\lambda_{X_s\cup X_{s'}}=\lambda_{X_{s'}}=:\lambda$ for all $s,s'>0$, the same restriction argument gives one $\sigma$ for all $X_s$, and $\kappa/\lambda$ fails (b) on $X_s$ for large $s$.
\end{proof}

\begin{remark}[Same minimisers instead of same values]\label{rem:argmin}
(a) If a fixed symmetric $\Delta$ satisfies $\arg\min J_{K+\Delta}=\arg\min J_K$ for every symmetric $K$, or only for every PD $K$, then $\Delta\in V_{n,k}$: otherwise take $K=\lambda I-\Delta$ with $\lambda$ large, so that $K$ and $K+\Delta=\lambda I$ are PD, $J_{K+\Delta}$ is constant (every partition is a minimiser) and $\arg\min J_K=\arg\max J_\Delta$ is a proper subset. Read as uniform transformations of the matrix, ``same minimisers'' adds nothing to Proposition~\ref{prop:invcomplete}. \status{proved here (one line)} (b) On a fixed dataset the weaker notion is vacuous: $K+\sigma I$ for $\sigma\ge\sigma_*$ is PSD with the same minimisers, and so is the block matrix $\sum_c\one_c\one_c^{\!\top}$ of any partition that is the unique minimiser. (c) What remains open is whether some dataset-independent PD kernel $\kappa$ has $\arg\min J_{\kappa[X]}=\arg\min J_{-A[X]}$ for every dataset and every $k$, for the rest-length spring. Corollary~\ref{cor:rest}(c) excludes the positive affine reparametrisations; a necessary condition from three-point sets with $k=2$ is that $J_K(\{i\},\{j,l\})=\tfrac12D(K)_{jl}$, so the pair $\{j,l\}$ that is closest in the feature metric of $\kappa$ must be the pair whose distance is closest to $r$. We do not know whether such a $\kappa$ exists, and the linear algebra used here does not reach ordinal conditions of this kind.
\end{remark}
%%%% BLOCK B END

%%%% BLOCK C BEGIN
\paragraph{Size weights.} The normalisation is the third route in the list above. Following \citet{dhillon2004,dhillon2007}, for point weights $\omega\in(0,\infty)^n$ fixed before the partition is chosen and a symmetric $K$, the \emph{weighted kernel $k$-means} objective is
\[
J^\omega_K(C)=\sum_i\omega_iK_{ii}-\sum_c\frac{\sum_{i,j\in C_c}\omega_i\omega_jK_{ij}}{\omega(C_c)},\qquad\omega(S)=\sum_{i\in S}\omega_i,
\]
which for $K=\Phi\Phi^{\!\top}$ is $\sum_c\sum_{i\in C_c}\omega_i\lVert\Phi_i-m_c\rVert^2$ with $m_c$ the $\omega$-weighted centroid \status{literature}. For a size weight $w$ write $F^w_A(C)=\sum_cw(n_c)\sum_{i<j\in C_c}A_{ij}$; only $w(m)$ for $2\le m\le n-k+1$ enters. Two facts are used. First, $\mathcal G_\omega=\{J^\omega_K:K\text{ symmetric}\}$ is a linear space of functions of the partition, of dimension at most $\binom n2$ (the $n$-dimensional space $\{g\one^{\!\top}+\one g^{\!\top}\}$ is in the kernel of $K\mapsto J^\omega_K$), and it contains the constants ($J^\omega_K\equiv n-k$ for $K=\mathrm{diag}(1/\omega_i)$); so ``up to a constant and a sign or positive factor'' adds nothing. Second, $F^w_A$ is linear in $A$ and depends only on its off-diagonal part.

\begin{proposition}[Size weights inside weighted kernel $k$-means, fixed point weights]\label{prop:weights}
Let $n\ge5$, $3\le k\le n-2$ and $w(2)\ne0$. The following are equivalent:
\begin{enumerate}[label=(\roman*),leftmargin=2em]
\item there is $\omega\in(0,\infty)^n$ such that for every symmetric $A$ there is a symmetric $K$ with $F^w_A-J^\omega_K$ constant on $k$-partitions;
\item the same for every Hooke spring matrix $A_{ij}=(x_i-x_j)^2$ of $n$ distinct points $x_i\in\R$;
\item $w(m)=2w(2)/m$ for $m=2,\dots,n-k+1$.
\end{enumerate}
Under (iii) one can take $\omega=\one$ and $K=-w(2)A$ (zero diagonal). For $k=2$ and $n\ge4$, (i) implies that $h(a)+h(n-a)$ does not depend on $a\in\{1,\dots,n-1\}$, where $h(m)=w(m)\binom m2$, provided $w(n-1)\ne0$ and $w(2)+w(n-2)\ne0$; for $n=4$ this is again (iii). \status{proved here}
\end{proposition}
\begin{proof}
(iii)$\Rightarrow$(i): with $w(m)=\lambda/m$, $F^w_A=\lambda E_\spec$ for the zero-diagonal $A$, and $E_\spec=\tfrac12J_{-A}$ by Theorem~\ref{thm:main}(b). (i)$\Rightarrow$(ii) is trivial. (ii)$\Rightarrow$(i): the matrices $M(x)_{ij}=(x_i-x_j)^2$ satisfy $M(x+y)-M(x)-M(y)=2\bigl[(x_i-x_j)(y_i-y_j)\bigr]_{ij}$, which for $x=e_a$, $y=e_b$ is $-2(E_{ab}+E_{ba})$; so they span all zero-diagonal symmetric matrices, and since $M$ is polynomial, so do the $M(x)$ with distinct coordinates (a non-empty open set). Membership in the linear space $\mathcal G_\omega$ passes to linear combinations. (i)$\Rightarrow$(iii): under (i), the $\binom n2$ functions $F^w_{E_{ab}+E_{ba}}$ ($a<b$) and the constant lie in $\mathcal G_\omega$, of dimension $\le\binom n2$, so they are linearly dependent: some $e\ne0$ (zero diagonal, symmetric) has $F^w_e$ constant. The swap of the proof of Proposition~\ref{prop:invcomplete}, with energies weighted by $w$, gives $w(2)(e_{pi}-e_{pj})=0$, so $e\equiv\varepsilon\ne0$ off the diagonal and $\sum_ch(n_c)$ is constant over all sizes $(n_1,\dots,n_k)$ of $k$-partitions. Compare the size vectors $(m+1,1,\rho)$ and $(m,2,\rho)$ with the same $k-2\ge1$ remaining sizes $\rho$ (possible for $1\le m\le n-k$): $h(m+1)-h(m)=h(2)-h(1)=h(2)$, so $h(m)=(m-1)h(2)$, i.e.\ $w(m)=2w(2)/m$. For $k=2$ the swap with $P=\emptyset$ and $P=\{x\}$ gives $w(n-1)\sum_{U}\gamma=0$ and $(w(2)+w(n-2))\gamma_x=0$, hence again $e\equiv\varepsilon$, and $\sum_ch(n_c)=h(a)+h(n-a)$; for $n=4$, $h(1)+h(3)=2h(2)$ reads $3w(3)=2w(2)$.
\end{proof}

\noindent Taking $n$ large, a size weight that is a weighted kernel $k$-means objective with fixed point weights for all $n$ and $k$ must be $w(m)\propto1/m$ for all $m\ge2$. The case $k=2$ with $h(a)+h(n-a)$ constant and $w\not\propto1/m$ is not settled by the dimension count. Point weights that depend on $A$, as in the normalised-cut equivalence of \citet{dhillon2004,dhillon2007}, are a weaker requirement; for the total energy they do not help:

\begin{proposition}[The total energy is not weighted kernel $k$-means]\label{prop:weightsdd}
Let $2\le k\le n-2$ and let $A$ be symmetric with $A_{ij}\ge0$ for $i\ne j$, not all zero (for instance the spring matrix of distinct points for any $\phi$ with $\phi(d)>0$ for $d>0$). Then for no $\omega\in(0,\infty)^n$, possibly depending on $A$, and no symmetric $K$ is $E_\tot-J^\omega_K$ constant on $k$-partitions. If $3\le k\le n-2$ or $(n,k)=(4,2)$, the same holds for $F^w_A$ whenever $w(3)>0$ and $w(3)\ge w(2)$. \status{proved here}
\end{proposition}
\begin{proof}
Take $a\ne b$ with $A_{ab}>0$. Split $[n]$ into $k+2$ non-empty groups with $a,b$ in different groups; keep $k-2$ groups $R_1,\dots,R_{k-2}$ as fixed clusters and call the other four $G_1,\dots,G_4$, with $a\in G_1$, $b\in G_2$; for the second statement use singletons $G_1=\{a\}$, $G_2=\{b\}$, $G_3=\{c\}$, $G_4=\{d\}$ (possible exactly when $k\ge3$ or $n=4$). For each of the seven partitions $\{X,Y\}$ of $\{1,2,3,4\}$ into two non-empty sets, let $C_X$ be the $k$-partition with clusters $G_X=\bigcup_{t\in X}G_t$, $G_Y$ and the $R$'s, and put
\[
\nu(C_X)=\varepsilon(X)\,\hat\omega(X)\,\hat\omega(Y),\qquad\hat\omega(X)=\omega(G_X),\quad\varepsilon(X)=\begin{cases}+1,&|X|\in\{1,3\},\\-1,&|X|=2,\end{cases}
\]
and $\nu=0$ on all other partitions. (1) $\sum_X\nu(C_X)=\bigl(\hat\Omega^2-\sum_t\hat\omega_t^2\bigr)-2\sum_{t<u}\hat\omega_t\hat\omega_u=0$, with $\hat\Omega=\sum_t\hat\omega_t$, because each pair $t<u$ is separated by exactly two of the three $2{+}2$ splits. (2) $\nu$ annihilates $\mathcal G_\omega$: by (1) the terms of $J^\omega_K(C_X)$ that are common to the seven partitions drop out, and what remains is $-\sum_{S\in\{X,Y\}}\hat Q(S)/\hat\omega(S)$ with $\hat Q(S)=\sum_{t,u\in S}\hat B_{tu}$, $\hat B_{tu}=\sum_{i\in G_t,j\in G_u}\omega_i\omega_jK_{ij}$. Every non-empty proper $S\subset\{1,2,3,4\}$ is a block of exactly one $C_X$, and $\nu(C_X)/\hat\omega(S)=\varepsilon(S)\hat\omega(S^c)$; so it suffices that $\sum_{S\ni t,u}\varepsilon(S)\hat\omega(S^c)=0$ for all $t,u$. For $t=u$: $(\hat\Omega-\hat\omega_t)-\sum_{v\ne t}(\hat\Omega-\hat\omega_t-\hat\omega_v)+\sum_{v\ne t}\hat\omega_v=0$; for $t\ne u$: $-\hat\omega(\{t,u\}^c)+\hat\omega(\{t,u\}^c)=0$. (3) For $E_\tot$, $E_\tot(C_X)=\mathrm{const}+\sum_{t<u,\ t,u\text{ on the same side}}\hat A_{tu}$ with $\hat A_{tu}=\sum_{i\in G_t,j\in G_u}A_{ij}\ge0$, and $\hat A_{12}\ge A_{ab}>0$. The pair $\{t,u\}$ lies in a two-group block once ($-\hat\omega(\{t,u\})\hat\omega(\{v,z\})$) and in a three-group block twice ($\hat\omega_v(\hat\Omega-\hat\omega_v)+\hat\omega_z(\hat\Omega-\hat\omega_z)$), where $\{v,z\}$ is the complementary pair; hence
\[
\langle\nu,E_\tot\rangle=2\sum_{t<u}\hat A_{tu}\,\hat\omega_v\hat\omega_z\ \ge\ 2A_{ab}\,\hat\omega_3\hat\omega_4>0 .
\]
Since $\langle\nu,J^\omega_K+\mathrm{const}\rangle=0$, $E_\tot$ is not of that form. With singleton groups and a size weight, the same count gives $\langle\nu,F^w_A\rangle=\sum_{t<u}A_{tu}\bigl[(w(3)-w(2))\hat\omega(\{t,u\})\hat\omega(\{v,z\})+2w(3)\hat\omega_v\hat\omega_z\bigr]>0$ under the stated conditions on $w$.
\end{proof}

\noindent For $w=1/m$ the pairing vanishes at $\omega=\one$, as it must. For decreasing weights with $w(3)<w(2)$ (e.g.\ $1/m^2$, $\binom m2^{-1}$) the bracket changes sign on $(0,\infty)^4$. So for $n=4$, $k=2$ a single spring can then be matched by suitable data-dependent weights, and Proposition~\ref{prop:weightsdd} says nothing about them; Proposition~\ref{prop:weights} still excludes fixed weights. Max-$k$-cut and min-sum $k$-clustering (Theorem~\ref{thm:main}(d)) are therefore not weighted kernel $k$-means objectives for non-negative spring matrices, unlike ratio association and normalised cut~\citep{dhillon2007}. We know of no statement of Propositions~\ref{prop:weights}--\ref{prop:weightsdd} in the literature, but have not searched exhaustively.
%%%% BLOCK C END

%%%% BLOCK D BEGIN
\begin{problem}[Dynamic mechanical clustering]\label{prob:dyn}
Let $V\colon(0,\infty)\to\R$ be a pair potential, let $z(t)$ solve $\dot z_i=-\sum_{j\ne i}V'(\lVert z_i-z_j\rVert)\,\frac{z_i-z_j}{\lVert z_i-z_j\rVert}$ with $z(0)=X$, and let $P_{T,\varepsilon}(X)$ be the single-linkage partition of $z(T)$ at threshold $\varepsilon$. Characterise the triples $(V,T,\varepsilon)$ for which there exist a fixed function $a\colon[0,\infty)\to\R$ and a normalisation $w$ such that, for every $X$, $P_{T,\varepsilon}(X)$ is the \emph{unique} minimiser of $\sum_cw(n_c)\sum_{i<j\in C_c}a(\lVert x_i-x_j\rVert)$ over partitions into the number of clusters it produces.
\end{problem}

\begin{remark}[Why ``unique'', and how a counterexample could be certified]\label{rem:dyn}
Without uniqueness the problem is trivial: for $a\equiv0$ every partition is a minimiser (and for $w=1/n_c$ so it is for every constant $a$). With uniqueness, a single configuration whose pairwise distances are all distinct is never a counterexample: choose $a=-1$ on the distances inside the blocks of $P_{T,\varepsilon}(X)$ and $a=0$ on the others; for $w\in\{1,1/n_c\}$ this makes $P_{T,\varepsilon}(X)$ the unique minimiser. A counterexample therefore needs repeated distances, within one configuration or across several configurations that share distances (e.g.\ subsets of a lattice). For finitely many configurations it is a finite system of strict linear inequalities in the values of $a$ on the common distance set, after scaling $F(C)-F(P)\ge1$ for $C\ne P$, so infeasibility can be certified exactly by a Farkas certificate over $\Q$. We have not carried this out.
\end{remark}
%%%% BLOCK D END
